\subsection{Principal findings}

The audit demonstrates that the public DGHS Dengue Dynamic Dashboard can serve as a bounded research-data interface when the requested state, returned response, displayed date, extracted fields, and integrity record are documented together. This is a reproducibility result for the frozen archive, not evidence that all dashboard history or all upstream reporting was captured. The empirical tables show why this distinction matters: raw labels, compatible universes, and typed longitudinal values cannot be assumed from visual availability alone.

The most consequential interpretation is that provenance and semantics are complementary. A response can be archived and hash-checked while its age category, geography, admission terminology, or reporting window remains incompletely defined. The study therefore makes an observation-to-table path inspectable without treating traceability as proof of population meaning. Source-labelled headline observations remain usable for restricted description, while normalized demographic or geographic comparisons remain limited by the contracts documented in B0.11.

\subsection{Relationship to existing Bangladesh dengue research}

Bangladesh already has national temporal and geographic analyses, district hotspot work, broad 64-district ecological analysis, surveillance-system assessment, and hospital/passive reporting literature \citep{ref1,ref2,ref3,ref4,ref6,ref7}. A publicly available non-peer-reviewed reproducible analysis repository also describes a DGHS HEOC dashboard-derived national analysis \citep{ref5}. ResearchDengu therefore does not add another burden, spatial, climate, predictive, or transmission analysis. It examines a different layer: the publicly exposed dashboard state as a research interface requiring provenance, semantic interpretation, and internal-consistency checks.

This distinction is substantive. Prior studies can use reported observations to answer epidemiological questions, while a dashboard audit asks whether a secondary researcher can reconstruct and interpret the exposed observation without silently changing its universe. An unresolved dashboard-level semantic issue does not establish a defect in the complete DGHS surveillance system, just as reproducible retrieval does not establish a complete national account. The present study complements, rather than replaces, epidemiological and system-level work.

The distinction also affects how apparent agreement should be read. A matching number across two panels may reflect shared display logic, but it does not demonstrate that the panels describe the same reporting population or time window. Conversely, a mismatch may be expected when one field is weekly and another cumulative, or when a geographic chart uses a different administrative universe. The reconciliation table therefore adds interpretive information even where it does not produce a corrected value: it tells a future user which comparisons are defensible and which should remain open.

\subsection{Implications for epidemiological researchers}

Acquisition should be treated as part of measurement documentation. Researchers should record the public submission, retrieval time, requested and displayed dates, response status, returned structure, raw response, and SHA-256 integrity hash before parsing. They should preserve source terminology and distinguish weekly, cumulative, daily, EPI-week, monthly, and annual fields until compatible definitions are established. A label that appears in a chart is not self-interpreting if its numerator, denominator, reporting pathway, or geographic assignment is undocumented.

Reconciliation should be attempted only after the compared fields are shown to share a reporting universe and time window. An unresolved or not-comparable result is preferable to forced totals, imputation, or a claim of system failure. Repeated captures are also required before a changed historical state can be characterized as a revision pattern. These practices allow researchers to use bounded dashboard observations without converting them into patient-residence, population, or community measures.

The sampling record has a related implication. Calendar anchors and pilot dates can make a corpus auditable, but they do not make it representative of every day or every reporting transition. Users should distinguish a predefined archive from a complete surveillance time series and should report whether a state was newly retrieved or reused from an earlier verified artifact. That distinction is especially important when a later analysis compares apparent changes across years or attempts to reproduce a historical display.

\subsection{Implications for public surveillance-data publication}

The findings suggest publication principles that could improve secondary use without imposing an obligation on DGHS. Stable metadata should identify case or admission definitions, reporting universe, time window, geographic meaning, and the relation between interval and cumulative fields. Machine-readable historical exports should retain release or retrieval dates, data dictionaries, stable field names, and version or change information. Geographic fields would be more interpretable with stable administrative codes and an explicit distinction between residence, treatment facility, reporting facility, and administrative aggregation. These are design recommendations, not claims that the current dashboard must implement them.

For secondary users, these principles would make an important distinction visible at the point of download: whether a value is a source observation, an updated version of a previous observation, or an analyst-derived summary. A documented change history would also allow a researcher to cite the state actually used rather than relying only on the date displayed in a browser. Clear labels for unavailable families and unresolved definitions would reduce pressure to fill gaps through ad hoc normalization. None of these recommendations changes the evidentiary status of the current corpus; they describe information that would make future public data states easier to reproduce and interpret.

\subsection{Strengths}

Strengths include a predefined acquisition/validation scope, archived response bodies with SHA-256 integrity hashes, raw-label preservation, semantic contracts, explicit reconciliation classifications, result-level provenance, independent numerical verification, and manual spot-validation. The study also separates restricted descriptive reporting from unsupported epidemiological inference. These features support auditability of the project corpus but do not remove the limitations below.

\subsection{Limitations}

The study evaluates publicly exposed dashboard states, not the complete underlying DGHS surveillance database. The dashboard may reflect hospital or passive surveillance; retrieval does not measure community ascertainment, health-seeking behavior, notification completeness, or all upstream records. Case, admission, and death semantics remain incompletely documented, and reporting-facility coverage may vary across dates. Geographic attribution is unresolved: district and City Corporation labels were not established as residence, treatment location, reporting facility, or another administrative unit.

Population denominators and numerator compatibility were not established, so the study does not estimate population incidence, prevalence, or individual risk. Historical coverage is limited to the validated snapshots and cadence points; 2026 is incomplete and is not a completed calendar year. The snapshot design cannot observe every intermediate dashboard state and does not establish continuous coverage. The 27-date selection rationale is recoverable from manifests at the level of calendar candidates and pilot inclusion, but the available records do not prove that all final choices preceded outcome viewing.

Ambiguous, malformed, overlapping, and spelling-variant categories were preserved rather than guessed. Reconciliation results remain classifications of compatibility limits, not corrected values. Revision behavior is limited to acquired snapshots: typed repeated series values were absent, so the study cannot characterize a prospective revision policy or demonstrate that revisions occurred. Summary-strip and last-24-hour values represent different windows and were not conflated. Finally, findings about this dashboard corpus should not be generalized automatically to all DGHS systems, all public dashboards, or Bangladesh's surveillance system as a whole.

The manuscript also does not establish that the observed category patterns arose from a particular software defect, reporting practice, clinical definition, or administrative decision. The audit records what was exposed and what could be compared under the frozen contracts. Resolving those mechanisms would require authoritative source documentation or additional evidence outside the current study and is therefore not treated as an R1 repair.
